\chapter{Configurația de lucru}
\label{ch:configurare}

În acest capitol descriem configurația concretă pe care a fost dezvoltat și testat proiectul: hardware, sistem de operare, alegerea de a nu folosi mașină virtuală, și versiunile exacte ale instrumentelor.

\section{Hardware}

Dezvoltarea s-a desfășurat în principal pe stația de lucru personală a șefei de echipă, cu următoarele specificații:

\begin{table}[!htb]
\centering
\caption{Configurația hardware folosită pentru dezvoltare și testare.}
\label{tab:hardware}
\begin{tabular}{ll}
\toprule
\textbf{Componentă} & \textbf{Specificație} \\
\midrule
Procesor & Intel Core i7-11800H, 8 nuclee, 16 fire, 2.30 GHz \\
Memorie RAM & 16 GB DDR4 \\
Stocare & SSD NVMe 512 GB \\
Placă video & NVIDIA GeForce RTX 3050 (4 GB) \\
\bottomrule
\end{tabular}
\end{table}

Cele trei membre ale echipei au folosit configurații similare, toate suficient de puternice pentru rularea suitei de teste și a instrumentelor de coverage. Suita completă (\texttt{mvn test}) rulează în aproximativ 2-3 secunde, iar PITest, care e cel mai costisitor, durează între 5 și 15 minute pentru cei 211 mutanți generați.

\section{Sistem de operare și mașini virtuale}

Sistemul de operare folosit a fost \textbf{Windows 11 Pro} (build 22H2), cu \textbf{PowerShell} ca shell pentru rularea comenzilor Maven. Comenzile sunt portabile și pe Linux/macOS dacă proiectul ar fi clonat acolo - Maven Wrapper (\texttt{mvnw}) e inclus în repository și are atât versiune \texttt{.cmd} pentru Windows, cât și script Unix.

\textbf{Nu a fost folosită nicio mașină virtuală.} Toate instrumentele (JDK, Maven, IntelliJ IDEA, Postgres pentru rularea aplicației) au fost instalate nativ pe Windows. Decizia a fost luată din considerente de performanță: rularea unei suite de teste cu Mockito e ușor afectată de overhead-ul unei VM, iar pentru PITest - care rulează \texttt{n} variante mutate ale codului, fiecare cu suita completă - timpul total ar fi crescut semnificativ.

Pentru baza de date am folosit \textbf{PostgreSQL 16} pentru aplicație și \textbf{H2 in-memory} pentru cazurile în care a fost nevoie să rulăm aplicația fără backend extern. Atenție însă: testele unitare nu folosesc nici Postgres, nici H2 - dependențele sunt complet mock-uite cu Mockito, deci suita rulează independent de orice bază de date.

\section{Stiva tehnologică}

Aplicația \textit{Călătorii} e construită pe Java 17 și Spring Boot 3.3.4. Tabelul \ref{tab:stack} sintetizează componentele majore.

\begin{table}[!htb]
\centering
\caption{Stiva tehnologică a aplicației și a suitei de teste.}
\label{tab:stack}
\begin{tabular}{lll}
\toprule
\textbf{Categorie} & \textbf{Tool / framework} & \textbf{Versiune} \\
\midrule
Limbaj & Java (OpenJDK) & 17 \\
Framework principal & Spring Boot & 3.3.4 \\
ORM & Hibernate (prin Spring Data JPA) & 6.5.x \\
Templating & Thymeleaf & 3.1.2 \\
Bază de date (prod) & PostgreSQL & 16 \\
Bază de date (dev) & H2 in-memory & 2.x \\
Build tool & Maven & 3.9.11 (prin wrapper) \\
Generare cod & Lombok & 1.18.42 \\
\midrule
Test framework & JUnit Jupiter & 5.10.3 \\
Mocking & Mockito (\texttt{mockito-junit-jupiter}) & 5.11.0 \\
Coverage & JaCoCo (Maven plugin) & 0.8.11 \\
Mutation Testing & PITest (\texttt{pitest-maven}) & 1.16.3 \\
\midrule
IDE & IntelliJ IDEA Ultimate & 2024.3.5 \\
Sistem de control versiuni & Git & 2.45.x \\
Hosting repo & GitHub & N/A \\
\bottomrule
\end{tabular}
\end{table}

\subsection{Justificarea alegerii instrumentelor de testare}

Stiva \textit{JUnit 5 + Mockito + JaCoCo + PITest} este, de facto, standardul pentru testare unitară pe ecosistemul Java în 2026. Câteva motive concrete pentru fiecare alegere:

\paragraph{JUnit 5 (Jupiter)} este versiunea modernă, modulară, a JUnit-ului. Față de JUnit 4 oferă: metode de test \texttt{package-private} (nu mai trebuie \texttt{public}), un sistem de extensii curat prin \texttt{@ExtendWith}, asserțiuni îmbunătățite (\texttt{assertThrows}, \texttt{assertAll}), suport nativ pentru parametrizate. Vine implicit prin \texttt{spring-boot-starter-test}, deci nu trebuie adăugat manual.

\paragraph{Mockito} este alegerea standard pentru mock-uri în Java, cu o sintaxă fluentă (\texttt{when(...).thenReturn(...)}) și integrare elegantă cu JUnit 5. Alternativele - EasyMock, PowerMock - sunt fie mai vechi, fie au probleme cu noile versiuni de Java.

\paragraph{JaCoCo} este, de asemenea, instrumentul standard în ecosistemul Java, integrat cu Maven, Gradle și majoritatea IDE-urilor. Generează rapoarte HTML interactive în care fiecare linie e colorată în verde (acoperită), galben (parțial acoperită) sau roșu (neatinsă).

\paragraph{PITest} este instrumentul lider pentru Mutation Testing pe Java, cu performanță bună și un set bogat de mutatori. Alternativele (\textit{MuJava}, \textit{Major}) sunt mai academice, mai puțin întreținute și mai dificil de integrat cu Maven.

\section{Arhitectura proiectului}

Aplicația respectă structura clasică \textit{layered} a unui proiect Spring Boot, vizibilă în figura \ref{fig:architecture}:

\begin{figure}[!htb]
    \centering
    \includegraphics[width=0.9\linewidth]{01-architecture.png}
    \caption{Arhitectura aplicației \textit{Călătorii} pe trei straturi (prezentare, business, persistență), cu evidențierea clasei testate \texttt{WalletServiceImpl}.}
    \label{fig:architecture}
\end{figure}

\begin{itemize}
    \item \textbf{Stratul de prezentare} - controllere Spring MVC + template-uri Thymeleaf. Pentru wallet, există un singur \texttt{WalletController} care expune endpoint-urile HTTP.
    \item \textbf{Stratul de business} - serviciile cu logica aplicativă. \texttt{WalletServiceImpl} e implementarea concretă a interfeței \texttt{WalletService} și e clasa pe care o testăm în acest proiect.
    \item \textbf{Stratul de persistență} - patru repository-uri JPA pe care \texttt{WalletServiceImpl} le folosește prin injecție de dependențe: \texttt{TripRepository}, \texttt{TripWalletRepository}, \texttt{UserRepository}, \texttt{WalletTransactionRepository}.
    \item \textbf{Baza de date} - PostgreSQL pentru rularea efectivă a aplicației.
\end{itemize}

În testele unitare, întreg stratul de persistență e mock-uit, ceea ce face ca testele să fie rapide și deterministe.

\section{Configurarea Maven (pom.xml)}

Plugin-urile principale relevante pentru testare sunt configurate în \texttt{pom.xml} astfel.

Pentru \textbf{JaCoCo}:

\begin{lstlisting}[language=XML, caption={Configurarea pluginului JaCoCo în pom.xml.}]
<plugin>
    <groupId>org.jacoco</groupId>
    <artifactId>jacoco-maven-plugin</artifactId>
    <version>0.8.11</version>
    <executions>
        <execution>
            <id>prepare-agent</id>
            <goals>
                <goal>prepare-agent</goal>
            </goals>
        </execution>
        <execution>
            <id>report</id>
            <phase>test</phase>
            <goals>
                <goal>report</goal>
            </goals>
            <configuration>
                <includes>
                    <include>echipa13/calatorii/service/impl/WalletServiceImpl.class</include>
                </includes>
            </configuration>
        </execution>
    </executions>
</plugin>
\end{lstlisting}

Filtrul \texttt{includes} limitează raportul JaCoCo la o singură clasă - \texttt{WalletServiceImpl}. Decizia e deliberată: cerința proiectului e să testăm o singură clasă reprezentativă, iar dacă raportul ar include toată aplicația, procentele ar fi diluate de codul netestat din celelalte module.

Pentru \textbf{PITest}:

\begin{lstlisting}[language=XML, caption={Configurarea pluginului PITest în pom.xml.}]
<plugin>
    <groupId>org.pitest</groupId>
    <artifactId>pitest-maven</artifactId>
    <version>1.16.3</version>
    <dependencies>
        <dependency>
            <groupId>org.pitest</groupId>
            <artifactId>pitest-junit5-plugin</artifactId>
            <version>1.2.2</version>
        </dependency>
    </dependencies>
    <configuration>
        <targetClasses>
            <param>echipa13.calatorii.service.impl.WalletServiceImpl</param>
        </targetClasses>
        <targetTests>
            <param>echipa13.calatorii.service.impl.WalletServiceImpl*Test</param>
        </targetTests>
        <outputFormats>
            <outputFormat>HTML</outputFormat>
            <outputFormat>XML</outputFormat>
        </outputFormats>
        <mutators>
            <mutator>STRONGER</mutator>
        </mutators>
    </configuration>
</plugin>
\end{lstlisting}

Setul de mutatori \texttt{STRONGER} extinde setul implicit (\texttt{DEFAULTS}) cu mutatori suplimentari: \textit{InlineConstantMutator}, \textit{NonVoidMethodCallMutator}, \textit{ExperimentalSwitchMutator} etc. Decizia de a folosi setul mai complet a fost luată pentru că vrem un raport care să stresseze cât mai mult suita.

Plugin-ul \texttt{pitest-junit5-plugin} e necesar deoarece varianta de bază a PITest înțelege doar JUnit 4. Trebuie declarat ca dependență internă a plugin-ului.

\section{Versiunea de Java și note de compatibilitate}

Proiectul rulează pe \textbf{Java 17} (LTS), care este versiunea minimă cerută de Spring Boot 3.x. Pe stațiile de dezvoltare, JDK-ul efectiv folosit a fost \textbf{OpenJDK 24}, care e backwards-compatible cu Java 17.

O incompatibilitate pe care am întâlnit-o în timpul setup-ului: versiunea standard de Lombok (1.18.32) nu funcționa cu OpenJDK 24, returnând erori de \textit{annotation processing}. Soluția a fost să forțăm explicit Lombok 1.18.42 (ultima versiune stabilă la momentul predării proiectului) prin \texttt{annotationProcessorPaths} în \texttt{maven-compiler-plugin}.

În capitolul următor descriem aplicația \textit{Călătorii} și clasa \texttt{WalletServiceImpl}, înainte de a intra în detaliile testării propriu-zise.
